\documentclass[../../main.tex]{subfiles}
\begin{document}

{\bf[解]}

% TODO 3D-tikzの円錐と円柱
\begin{figure}[htb]
  \centering
  % 手書きスケッチのアングル（俯角 65°，方位角 115°）
  \tdplotsetmaincoords{65}{115}
  \begin{tikzpicture}[tdplot_main_coords, scale=2.6, line join=round, >=stealth, every node/.style={font=\small}]

    % --- 1. 定数・座標パラメータ ---
    \coordinate (O)     at (0, 0, 0);
    \coordinate (V_top) at (0, 0, 1.0);  % 円錐の頂点 (0, 0, 1)
    \coordinate (C_cyl) at (1.0, 0, 0);  % 円柱底面の中心 (1, 0, 0)
    \coordinate (X2)    at (2.0, 0, 0);  % (2, 0, 0)

    % --- 2. 座標軸 ---
    \draw[dashed, gray!40] (-0.8, 0, 0) -- (0, 0, 0);
    \draw[dashed, gray!40] (0, -0.8, 0) -- (0, 0, 0);
    \draw[->, thick, gray!80!black] (0, 0, 0) -- (2.6, 0, 0) node[below left=1pt] {$x$};
    \draw[->, thick, gray!80!black] (0, 0, 0) -- (0, 2.4, 0) node[right=2pt] {$y$};
    \draw[->, thick, gray!80!black] (0, 0, 0) -- (0, 0, 1.6) node[above=2pt] {$z$};
    \node[below left=2pt] at (O) {$O$};

    % --- 3. 円柱側面全体の薄い斜線（薄い母線・背景ハッチング）---
    % 手前側および外側全体に薄い縦縞を描画 (z=0 から z=1 まで)
    \foreach \ang in {0, 15, ..., 345} {
      \draw[gray!35, thin]
      ({1.0 + cos(\ang)}, {sin(\ang)}, 0) --
      ({1.0 + cos(\ang)}, {sin(\ang)}, 1.0);
    }

    % --- 4. 円柱側面のうち「円錐の内部」にある部分の強調（濃い青の斜線・塗り）---
    % 円柱面: x = 2*cos^2(phi), y = 2*cos(phi)*sin(phi)
    % 円錐の高さ上限: z = 1 - cos(phi)  (-90° <= phi <= 90°)
    %
    % (i) 背景のカラーフィル
    \fill[blue!18, opacity=0.75]
    plot[domain=-90:90, samples=60] ({2.0*cos(\x)*cos(\x)}, {2.0*cos(\x)*sin(\x)}, {1.0 - cos(\x)})
    -- plot[domain=90:-90, samples=60] ({2.0*cos(\x)*cos(\x)}, {2.0*cos(\x)*sin(\x)}, 0)
    -- cycle;

    % (ii) 手書き風の密な斜線・縦縞（底面から交線への母線スライス）
    \foreach \phi in {-85, -75, ..., 85} {
      \draw[blue!85!black, line width=1.1pt]
      ({2.0*cos(\phi)*cos(\phi)}, {2.0*cos(\phi)*sin(\phi)}, 0) --
      ({2.0*cos(\phi)*cos(\phi)}, {2.0*cos(\phi)*sin(\phi)}, {1.0 - cos(\phi)});
    }

    % (iii) 交差曲線（円柱側面と円錐側面の交線: 太線実線）
    \draw[line width=1.8pt, blue!90!black, domain=-90:90, samples=70]
    plot ({2.0*cos(\x)*cos(\x)}, {2.0*cos(\x)*sin(\x)}, {1.0 - cos(\x)});

    % --- 5. 底面 z = 0 の輪郭線 ---
    % 円柱底面の輪郭
    \draw[dashed, gray!60, thick]
    plot[domain=0:360, samples=60] ({1.0 + cos(\x)}, {sin(\x)}, 0);

    % 円錐底面の輪郭 (手前 y >= 0 を実線、奥を破線)
    \draw[very thick, gray!75!black, domain=0:360, samples=60]
    plot ({2.0*cos(\x)}, {2.0*sin(\x)}, 0);

    % --- 6. 円柱 B の上面 (z = 1) と側面外枠 ---
    % 上面円 (z = 1, 中心 (1, 0, 1), 半径 1)
    \draw[thick, gray!75!black]
    plot[domain=0:360, samples=60] ({1.0 + cos(\x)}, {sin(\x)}, 1.0);
    \draw[dashed, gray!50] (1.0, 0, 1.0) -- (2.0, 0, 1.0);
    \draw[dashed, gray!50] (1.0, 0, 0) -- (1.0, 0, 1.0);

    % 円柱の最外殻母線 (x = 2)
    \draw[thick, gray!75!black] (2.0, 0, 0) -- (2.0, 0, 1.0);

    % --- 7. 円錐 A の母線 ---
    \draw[very thick, black] (V_top) -- ({2.0*cos(0)}, {2.0*sin(0)}, 0);   % (2, 0, 0)
    \draw[very thick, black] (V_top) -- ({2.0*cos(90)}, {2.0*sin(90)}, 0); % (0, 2, 0)
    \draw[very thick, black] (V_top) -- ({2.0*cos(180)}, {2.0*sin(180)}, 0); % (-2, 0, 0)

    % --- 8. 目盛りと主要点プロット ---
    \fill[black] (V_top) circle (1.0pt) node[left=2pt] {$1$};
    \fill[black] (C_cyl) circle (0.8pt) node[below=2pt] {$1$};
    \fill[black] (X2)    circle (0.8pt) node[below=2pt] {$2$};

  \end{tikzpicture}
  \caption{円錐 $A$ と円柱 $B$}
  \label{fig:cylinder_side_inside_cone_3d_shaded}
\end{figure}

円錐$A$，円柱$B$の方程式は
\begin{align}
  \begin{cases}
    A: x^2+y^2     \le (1-z)^2 \\
    B: (x-1)^2+y^2 \le 1
  \end{cases} \label{eq:1}
\end{align}
である．以下，簡単のため$c=\cos\theta, s=\sin\theta$とおく．

\paragraph{(1)}

\cref{eq:1}の$z=1-c$での切り口は，
\begin{align}
  \begin{cases}
    x^2+y^2 \le (2c)^2 \\
    (x-1)^2+y^2 \le 1
  \end{cases} \label{eq:2}
\end{align}
である．これを図示すると下図の斜線部となる．

% TODO: 切断面
\begin{figure}[htb]
  \centering
  \begin{tikzpicture}[scale=2.6, >=stealth, every node/.style={font=\small}]

    % --- 1. 定数・厳密幾何パラメータ ---
    % theta = 55°
    \def\angTheta{55}
    \pgfmathsetmacro{\cVal}{cos(\angTheta)}
    \pgfmathsetmacro{\sVal}{sin(\angTheta)}
    \pgfmathsetmacro{\rCone}{2.0 * \cVal}           % ≈ 1.1472
    \pgfmathsetmacro{\xP}{2.0 * \cVal * \cVal}      % 2c^2 ≈ 0.6580
    \pgfmathsetmacro{\yP}{2.0 * \cVal * \sVal}      % 2cs ≈ 0.9397

    \coordinate (O)     at (0, 0);
    \coordinate (A)     at (1.0, 0);                 % 円柱の中心 (1, 0)
    \coordinate (P_pos) at (\xP, \yP);
    \coordinate (P_neg) at (\xP, -\yP);
    \coordinate (Q_mid) at (\xP, 0);

    % --- 2. 切断面の共通部分のハッチング塗りつぶし ---
    \fill[pattern=north east lines, pattern color=blue!65, opacity=0.75]
    plot[domain=-\angTheta:\angTheta, samples=60] ({\rCone*cos(\x)}, {\rCone*sin(\x)})
    -- plot[domain={180-(180-2*\angTheta)}:{180+(180-2*\angTheta)}, samples=60] ({1.0 + cos(\x)}, {sin(\x)})
    -- cycle;

    % --- 3. 2つの円の輪郭線 ---
    % 円錐の断面円: x^2 + y^2 = (2c)^2
    \draw[thick, gray!50] (O) circle (\rCone);
    % 円柱の断面円: (x - 1)^2 + y^2 = 1
    \draw[thick, gray!50] (A) circle (1.0);

    % 共通部分の境界線（太線強調）
    \draw[very thick, blue!85!black]
    plot[domain=-\angTheta:\angTheta, samples=60] ({\rCone*cos(\x)}, {\rCone*sin(\x)});
    \draw[very thick, teal!85!black]
    plot[domain={180-(180-2*\angTheta)}:{180+(180-2*\angTheta)}, samples=60] ({1.0 + cos(\x)}, {sin(\x)});

    % --- 4. 座標軸 ---
    \draw[->, thick, gray!80!black] (-1.35, 0) -- (2.3, 0) node[right] {$x$};
    \draw[->, thick, gray!80!black] (0, -1.35) -- (0, 1.45) node[above] {$y$};
    \node[below left=2pt] at (O) {$O$};

    % --- 5. 補助線と直角・角度 ---
    % 弦 P_+ P_-
    \draw[dashed, gray!60, thick] (P_pos) -- (P_neg);

    % 動径 OP_+ (長さ 2c)
    \draw[very thick, blue!80!black] (O) -- (P_pos);
    \node[blue!90!black, font=\scriptsize, above left] at ($ (O)!0.5!(P_pos) $) {$2c$};

    % 動径 AP_+ (長さ 1)
    \draw[very thick, teal!80!black] (A) -- (P_pos);
    \node[teal!90!black, font=\scriptsize, above right] at ($ (A)!0.5!(P_pos) $) {$1$};

    % 角度 theta = angle(P_+ O Q)
    \draw[thick, purple!85!black] (0.35, 0) arc (0:\angTheta:0.35);
    \node[purple!90!black, font=\scriptsize\bfseries] at ({\angTheta/2}:0.48) {$\theta$};

    % 角度 pi - 2theta = angle(P_+ A Q)
    % A(1, 0) から見た P_+ の偏角は 180° - (180° - 2theta) = 2theta ではなく、
    % A から Q への方向は 180°。P_+ への方向は 180° - (180° - 2theta) = 180° - 70° = 110°
    \pgfmathsetmacro{\angAP}{180 - (180 - 2*\angTheta)} % 110°
    \draw[thick, orange!90!black] ($ (A) + (180:0.3) $) arc (180:\angAP:0.3);
    \node[orange!95!black, font=\tiny\bfseries] at ($ (A) + ({180 - (180-2*\angTheta)/2}:0.5) $) {$\pi-2\theta$};

    % --- 6. 頂点プロットとラベル ---
    \fill[black] (O)     circle (0.8pt);
    \fill[black] (A)     circle (0.8pt) node[below=2pt] {$A(1,0)$};
    \fill[red!85!black]  (P_pos) circle (1.2pt) node[above right=1pt, red!95!black] {$P_+(2c^2, 2cs)$};
    \fill[red!85!black]  (P_neg) circle (1.2pt) node[below right=1pt, red!95!black] {$P_-(2c^2, -2cs)$};
    \fill[gray!80!black] (Q_mid) circle (0.8pt) node[below right=-1pt, font=\scriptsize] {$Q$};

    % 目盛り
    \node[below=2pt, font=\scriptsize] at (\rCone, 0) {$2c$};
    \node[below=2pt, font=\scriptsize] at (2.0, 0) {$2$};

  \end{tikzpicture}
  \caption{切断面 $z = 1-c$ における共通領域（斜線部）と交点 $P_\pm$}
  \label{fig:cross_section_common_region}
\end{figure}

ここで，断面の$2$円の交点の$x$座標は\cref{eq:2}で不等号を等号にしたものを辺々引いて
\begin{align}
  x^2 - (x-1)^2 = (2c)^2-1 \quad \therefore x = 2c^2
\end{align}
である．この時$y$座標は
\begin{align}
  y^2 = (2c)^2 - (2c^2)^2 = 4c^2(1-c^2) = 4c^2s^2 \quad \therefore y = \pm 2cs
\end{align}
である．従って交点の座標は$P_{\pm}(2c^2,\pm 2cs)$である．
よって図において
\begin{align}
  \angle P_{+}OQ = \theta, \quad \angle P_{+}AQ = \pi - 2\theta
\end{align}
である．

\begin{figure}[htb]
  \centering
  \begin{tikzpicture}[scale=2.2, >=stealth, every node/.style={font=\small}]

    % ============================================================
    % 【左側】S_1 の半分 (中心 O, 半径 2c, 中心角 theta)
    % ============================================================
    \begin{scope}[xshift=0cm]
      \def\angT{48}
      \def\radOne{1.8} % 2c
      \pgfmathsetmacro{\xChord}{\radOne * cos(\angT)}

      \coordinate (O1) at (0, 0);

      % 弓形部分の斜線ハッチング (垂直線から右端の円弧まで)
      \fill[pattern=north east lines, pattern color=blue!85, opacity=0.75]
      (\xChord, 0)
      -- (\xChord, {\radOne * sin(\angT)})
      arc (\angT:0:\radOne)
      -- cycle;

      % 扇形の輪郭線
      \draw[very thick, black] (O1) -- (\radOne, 0);
      \draw[very thick, black] (O1) -- ({\angT}:\radOne);
      \draw[very thick, black] (\radOne, 0) arc (0:\angT:\radOne);

      % 垂直な弦（破線）
      \draw[dashed, very thick, black] (\xChord, 0) -- (\xChord, {\radOne * sin(\angT)});

      % 直角記号
      \draw[thin, black]
      ($ (\xChord, 0) + (-0.08, 0) $) --
      ($ (\xChord, 0) + (-0.08, 0.08) $) --
      ($ (\xChord, 0) + (0, 0.08) $);

      % 角度 theta の弧
      \draw[thick, black] (0.38, 0) arc (0:\angT:0.38);
      \node[font=\footnotesize\bfseries] at ({\angT/2}:0.52) {$\theta$};

      % 半径ラベル 2c
      \node[below=3pt, font=\footnotesize\bfseries] at (\radOne, 0) {$2c$};

      % 見出し
      \node[below=14pt, font=\footnotesize\bfseries] at ({\radOne/2}, 0) {$S_1$ の半分};
    \end{scope}

    % ============================================================
    % 【右側】S_2 の半分 (中心 A が右側, 半径 1, 中心角 2*theta)
    % ============================================================
    \begin{scope}[xshift=3.8cm]
      % 手書きスケッチ通り、頂点(中心)を右側に配置
      % 頂角 2*theta = 60°
      \def\angTwoT{60}
      \def\radTwo{1.6} % 半径 1 に対応
      \pgfmathsetmacro{\wBase}{\radTwo * cos(\angTwoT)}
      \pgfmathsetmacro{\hTop}{\radTwo * sin(\angTwoT)}

      % 中心 A を原点 (wBase, 0) に置く
      \coordinate (A_center) at (\wBase, 0);
      \coordinate (Top_pt)   at (0, \hTop);
      \coordinate (Foot_pt)  at (0, 0);
      % 左端の円弧端点: A_center から半径 radTwo で 180° 方向
      \coordinate (Left_pt)  at ({\wBase - \radTwo}, 0);

      % 弓形部分の斜線ハッチング (左側の円弧から垂直線 Foot_pt -- Top_pt まで)
      \fill[pattern=north east lines, pattern color=teal!85, opacity=0.75]
      (Foot_pt) -- (Top_pt)
      arc ({180 - \angTwoT}:180:\radTwo)
      -- cycle;

      % 扇形の輪郭線
      \draw[very thick, black] (Left_pt) -- (A_center); % 底辺
      \draw[very thick, black] (A_center) -- (Top_pt);  % 斜辺
      % 正しい円弧: 中心 A_center からの円弧 (Top_pt から Left_pt へ)
      \draw[very thick, black] (Top_pt) arc ({180 - \angTwoT}:180:\radTwo);

      % 垂直な弦（破線）
      \draw[dashed, very thick, black] (Foot_pt) -- (Top_pt);

      % 直角記号
      \draw[thin, black]
      ($ (Foot_pt) + (0.08, 0) $) --
      ($ (Foot_pt) + (0.08, 0.08) $) --
      ($ (Foot_pt) + (0, 0.08) $);

      % 角度 2*theta の弧と引き出し線
      \draw[thick, black] ($ (A_center) + (180:0.35) $) arc (180:{180 - \angTwoT}:0.35);
      \draw[thin, black] ($ (A_center) + ({180 - \angTwoT/2}:0.35) $) to[out=45, in=180] ++(0.3, 0.15)
      node[right=1pt, font=\footnotesize\bfseries] {$2\theta$};

      % 半径ラベル 1
      \node[below=3pt, font=\footnotesize\bfseries] at (Left_pt) {$1$};

      % 見出し
      \node[below=14pt, font=\footnotesize\bfseries] at ({(\wBase - \radTwo + \wBase)/2}, 0) {$S_2$ の半分};
    \end{scope}

  \end{tikzpicture}
  \caption{扇形から直角三角形を引いて $S_1, S_2$ を求める様子（手書きスケッチの配置）}
  \label{fig:two_circular_segments_sketch_side_by_side}
\end{figure}

この面積 $S(t)$ は上図の斜線部の面積 $S_1$ と$S_2$の合計である．すなわち，
\begin{align}
  S(t)=S_1+S_2 \label{eq:3}
\end{align}
である．

それぞれ扇形から対応する三角形の面積を引いたもので，各々計算すると，
\begin{align}
  S_1
  &=\dfrac{1}{2}2\theta(2c)^2 - \dfrac{1}{2}(2c^2)(cs)
  &=4c^2\theta-4c^3s
\end{align}
および
\begin{align}
  S_2
  &=\dfrac{1}{2}(2\pi-4\theta)-\dfrac{1}{2}2\sin(\pi-2\theta)\cos(\pi-2\theta) \\
  &=\pi-2\theta+2cs(1+2c^2) \\
  &=\pi-2\theta+2cs+4c^3s
\end{align}
だから\cref{eq:3}に代入して
\begin{align}
  S(t)
  &=(4c^2-2)\theta-4c^3s+\pi+2cs+4c^3s \nonumber\\
  &=2\theta(1-2s^2)+\pi-2cs \label{eq:4}
\end{align}
である．$\cdots$(答)

\paragraph{(2)}

\cref{eq:4}から求める体積$V$として，
\begin{align}
  V&=\int_0^1S(t)dt \\
  &=\int_0^{\pi/2}S(t)\dfrac{dt}{d\theta}d\theta \quad (t=1-c)\\
  &=\int_0^{\pi/2}S(t)sd\theta \\
  &=\int_0^{\pi/2}(2s\theta-4s^3\theta+\pi s-2cs^2)d\theta \quad \left(\because\text{\cref{eq:4}}\right) \label{eq:5}
\end{align}
である．各項計算すると
\begin{align}
  (\text{第一項})
  &= \int_0^{\pi/2}s\theta d\theta=2\left[-c\theta+s\right]_0^{\pi/2}=1
\end{align}
\begin{align}
  (\text{第二項})
  &= 4\int_0^{\pi/2}s^3\theta d\theta=\int_0^{\pi/2}(3s-\sin3\theta)\theta d\theta \\
  &=3-\left[\dfrac{-1}{3}\theta\cos3\theta+\dfrac{1}{9}\sin3\theta\right]_0^{\pi/2} \\
  &=3+\dfrac{1}{9}=\dfrac{28}{9}
\end{align}
\begin{align}
  (\text{第三項})
  &= \int_0^{\pi/2}sd\theta=1
\end{align}
\begin{align}
  (\text{第四項})
  &=\int_0^{\pi/2}cs^2d\theta=\left[\dfrac{1}{3}s^3\right]_0^{\pi/2}=\dfrac{1}{3}
\end{align}
だから，これらを\cref{eq:5}に代入して求める体制は
\begin{align}
  V=2-\dfrac{28}{9}+\pi-\dfrac{2}{3}=\pi-\dfrac{16}{9}
\end{align}
である．$\cdots$(の答)

{\bf [別解]}

\paragraph{(1)}

%TODO: upgrade
\begin{figure}[htb]
  \centering
  \begin{tikzpicture}[scale=2.6, >=stealth, every node/.style={font=\small}]

    % --- 1. 定数・厳密パラメータ (theta = 55°) ---
    \def\angTheta{55}
    \pgfmathsetmacro{\cVal}{cos(\angTheta)}
    \pgfmathsetmacro{\sVal}{sin(\angTheta)}
    \pgfmathsetmacro{\rCone}{2.0 * \cVal}
    \pgfmathsetmacro{\xP}{2.0 * \cVal * \cVal}
    \pgfmathsetmacro{\yP}{2.0 * \cVal * \sVal}

    \coordinate (O)     at (0, 0);
    \coordinate (A)     at (1.0, 0);
    \coordinate (P_pos) at (\xP, \yP);
    \coordinate (P_neg) at (\xP, -\yP);

    % --- 2. 積分領域 S1, S2 の塗りつぶし ---
    % S2 (左側): x in [1 - 2c^2, 1] relative to circle centered at A
    \fill[pattern=north west lines, pattern color=teal!70, opacity=0.8]
    (P_pos) -- plot[domain={180-(180-2*\angTheta)}:{180+(180-2*\angTheta)}, samples=50] ({1.0 + cos(\x)}, {sin(\x)})
    -- (P_neg) -- cycle;

    % S1 (右側): x in [2c^2, 2c]
    \fill[pattern=north east lines, pattern color=blue!70, opacity=0.8]
    (P_pos) -- plot[domain=\angTheta:-\angTheta, samples=50] ({\rCone*cos(\x)}, {\rCone*sin(\x)})
    -- (P_neg) -- cycle;

    % --- 3. 円周の描画 ---
    \draw[thick, gray!55] (O) circle (\rCone);
    \draw[thick, gray!55] (A) circle (1.0);

    % 共通境界の太線強調
    \draw[very thick, blue!85!black]
    plot[domain=\angTheta:-\angTheta, samples=50] ({\rCone*cos(\x)}, {\rCone*sin(\x)});
    \draw[very thick, teal!85!black]
    plot[domain={180-(180-2*\angTheta)}:{180+(180-2*\angTheta)}, samples=50] ({1.0 + cos(\x)}, {sin(\x)});

    % 分割境界 x = 2c^2
    \draw[very thick, red!85!black] (P_pos) -- (P_neg);

    % --- 4. 座標軸 ---
    \draw[->, thick, gray!80!black] (-1.35, 0) -- (2.3, 0) node[right] {$x$};
    \draw[->, thick, gray!80!black] (0, -1.35) -- (0, 1.45) node[above] {$y$};
    \node[below left=2pt] at (O) {$O$};

    % --- 5. 角度 theta の動径 ---
    \draw[thick, blue!80!black] (O) -- (P_pos);
    \draw[thick, purple!85!black] (0.35, 0) arc (0:\angTheta:0.35);
    \node[purple!90!black, font=\scriptsize\bfseries] at ({\angTheta/2}:0.48) {$\theta$};

    % --- 6. 積分境界の目盛りと点プロット ---
    \fill[black] (P_pos) circle (1.1pt) node[above right=1pt] {$P_+$};
    \fill[black] (P_neg) circle (1.1pt) node[below right=1pt] {$P_-$};
    \fill[black] (A)     circle (0.8pt) node[below=2pt] {$A$};

    % x 軸上の境界値
    \node[below=2pt, font=\scriptsize, red!90!black] at (\xP, 0) {$2c^2$};
    \node[below=2pt, font=\scriptsize] at (\rCone, 0) {$2c$};
    \node[below=2pt, font=\scriptsize] at (2.0, 0) {$2$};

    % 領域ラベル
    \node[blue!90!black, font=\footnotesize\bfseries] at (0.85, -0.3) {$S_1$};
    \node[teal!90!black, font=\footnotesize\bfseries] at (0.35, -0.3) {$S_2$};

  \end{tikzpicture}
  \caption{定積分で求める 2 つの領域 $S_1, S_2$ の積分区間と幾何関係}
  \label{fig:cross_section_integration_upgrade}
\end{figure}

$S(t)$の求積には通常通りの積分を用いても良い．
[解答]の中で定義した$S_1,S_2$は
\begin{align}
  S_1
  &=2\int_{2c^2}^{2c}\sqrt{(2c)^2-x^2}\,dx,\\
  S_2
  &=2\int_{1-2c^2}^{1}\sqrt{1-x^2}\,dx.
\end{align}
だから，それぞれ計算すると
\begin{align}
  \begin{split}
    \dfrac{1}{2}S_1
    &=(2c)^2\int_0^\theta\sin^2u\,du\\
    &=2c^2\left[u-\dfrac{1}{2}\sin2u\right]_0^\theta\\
    &=2c^2\left(\theta-\dfrac{1}{2}\sin2\theta\right),
  \end{split}\\
  \begin{split}
    \dfrac{1}{2}S_2
    &=\int_0^{\pi-2\theta}\sin^2u\,du\\
    &=\dfrac{1}{2}\left[u-\dfrac{1}{2}\sin2u\right]_0^{\pi-2\theta}\\
    &=\dfrac{1}{2}\left(\pi-2\theta+\dfrac{1}{2}\sin4\theta\right).
  \end{split}
\end{align}
だから
\begin{align}
  \begin{split}
    S(t)=S_1+S_2
    &=2\left[(2c^2-1)\theta+\dfrac{\pi}{2}
    -c^2\sin2\theta+\dfrac{1}{4}\sin4\theta\right]\\
    &=2\theta\cos2\theta+\pi-\sin2\theta.
  \end{split}
\end{align}
である．

\newpage
\end{document}
